\section{项目时间安排}
\label{sec:schedule}
项目按“数据审计、Caption 训练与指标、任务二迁移、Qwen 配对实验、SCST 与 Judge”推进，每阶段完成可运行入口及验收记录后进入下一阶段。图~\ref{fig:team-timeline}给出初始排期；模型接口开发、文档与离线复核并行，正式 GPU 作业按依赖顺序执行。当前已有 CPU tiny 链路，正式训练入口和 GPU 效果仍待验证，以下日期属于计划检查点。

\begin{figure}[!ht]
\centering% Weekly task windows; final evaluation dependencies are expanded below.
\begin{tikzpicture}[x=1.24cm,y=-1cm,font=\fontsize{8}{10}\selectfont]
\foreach \x/\lab in {0/11.04,1/11.09,2/11.16,3/11.23,4/11.30,5/12.07,6/12.14,7/12.21,8/12.28,9/01.04,10/01.11}{
 \node[text=FLBlue,font=\fontsize{7.5}{9}\selectfont] at (\x+.5,0) {\lab};
 \draw[FLLine!45] (\x,.23) -- (\x,4.7);}
\draw[FLLine!45] (11,.23)--(11,4.7);
\foreach \y/\lab/\clr in {.8/罗开炜/FLBlue,1.85/郝天乐/FLTeal,2.9/张嘉荣/FLAmber,3.95/GPU/FLMid}{
 \fill[\clr!4] (0,\y-.44) rectangle (11,\y+.44);
 \node[anchor=east,text=\clr,font=\small\bfseries] at (-.12,\y){\lab};}
\def\slot#1#2#3#4#5{\filldraw[rounded corners=1.7pt,draw=#4!70,fill=#4!10] (#1+.035,#3-.34) rectangle (#2-.035,#3+.34);
\node[text=FLText,align=center] at ({(#1+#2)/2},#3){#5};}
\slot{0}{2}{.8}{FLBlue}{方案与接口}
\slot{2}{4.45}{.8}{FLBlue}{训练器与模型联调}
\slot{4.45}{6.43}{.8}{FLBlue}{迁移开发\\适配复核}
\slot{6.43}{7}{.8}{FLBlue}{\fontsize{7}{9}\selectfont SCST}
\slot{7}{8}{.8}{FLBlue}{分歧\\复核}
\slot{8}{11}{.8}{FLBlue}{复现检查与技术整合}
\slot{0}{3}{1.85}{FLTeal}{数据审计、指标与试跑}
\slot{3}{5.57}{1.85}{FLTeal}{Caption 与迁移评测}
\slot{5.57}{7}{1.85}{FLTeal}{Qwen\\统一评测}
\slot{7}{8}{1.85}{FLTeal}{盲评\\复算}
\slot{8}{11}{1.85}{FLTeal}{误差分层与独立复算}
\slot{0}{2}{2.9}{FLAmber}{调研与开题}
\slot{2}{5}{2.9}{FLAmber}{分析工具、案例与中期材料}
\slot{5}{7}{2.9}{FLAmber}{Judge 接入\\人评校准}
\slot{7}{8}{2.9}{FLAmber}{盲评\\统计}
\slot{8}{11}{2.9}{FLAmber}{统计分析、展示与报告}
\slot{2.4}{3}{3.95}{FLMid}{\fontsize{7}{9}\selectfont 试跑}
\slot{3}{4.45}{3.95}{FLMid}{A/B\\含消融}
\slot{4.45}{5.57}{3.95}{FLMid}{Hybrid}
\slot{5.57}{6.43}{3.95}{FLMid}{Qwen}
\slot{6.43}{7}{3.95}{FLMid}{\fontsize{7}{9}\selectfont SCST\\\fontsize{7}{9}\selectfont +CE}
\slot{7}{8}{3.95}{FLMid}{测试生成\\Judge}
\node[flnote,text width=3.4cm,font=\fontsize{7.5}{9}\selectfont] at (9.5,3.95)
 {冻结配置下按需复跑\\离线分析不持续占用 GPU};
\foreach \x/\lab in {2/开题,5/中期,7/冻结,11/结题}{
 \fill[FLBlue] (\x,4.54)--(\x+.065,4.63)--(\x,4.72)--(\x-.065,4.63)--cycle;
 \node[anchor=north,text=FLBlue,font=\fontsize{7.5}{9}\selectfont] at (\x,4.75) {\lab};}

\node[fllabel,anchor=west] at (-1.12,5.5) {正式评价与汇报的先后关系};
\node[flnode,text width=2.55cm,minimum height=1.05cm] (freeze) at (.12,6.4)
 {\textbf{12/20}\\冻结模型与协议};
\node[flaccent,text width=2.55cm,minimum height=1.05cm] (generate) at (2.55,6.4)
 {\textbf{12/21 起}\\测试生成、Judge};
\node[flaccent,text width=2.55cm,minimum height=1.05cm] (human) at (4.98,6.4)
 {\textbf{12/22--24}\\双人独立盲评};
\node[fltrain,text width=2.55cm,minimum height=1.05cm] (merge) at (7.41,6.4)
 {\textbf{12/25}\\统计复核与合并};
\node[flnode,text width=2.55cm,minimum height=1.05cm] (present) at (9.84,6.4)
 {\textbf{12/26--28}\\演练与预留汇报};
\draw[flarrow] (freeze.east)--(generate.west);
\draw[flarrow] (generate.east)--(human.west);
\draw[flarrow] (human.east)--(merge.west);
\draw[flarrow] (merge.east)--(present.west);
\node[flnote,anchor=west,font=\fontsize{8}{10}\selectfont] at (-1.12,7.2)
 {盲评仅使用匿名候选；Judge 结果对评分者隐藏。图示为初始排期，依 GPU 试跑修订。};
\end{tikzpicture}


\caption{团队工作、单卡窗口与正式评价衔接。横轴为各周起始日期，1 月为 2027 年；菱形标记检查点。下方展开冻结、生成、盲评及合并顺序，GPU 行表示作业窗口而非持续占用。}
\label{fig:team-timeline}
\end{figure}

\textbf{11 月 4--22 日：方案、数据与运行条件就绪。}11 月 15 日完成开题提交；22 日前由郝天乐固定 Food500 图文映射、重复组、排除清单与划分，同时核验任务二原始数据。罗开炜实现自主训练器，郝天乐接入四指标，张嘉荣梳理文献依据并准备分析 Notebook。利用少量 GPU 时间预检查 Qwen 加载、监督掩码和适配器保存，检查 Judge 解析及 SCST 奖励接口，不将接口试跑作为正式实验结果。首次试跑记录批量、精度、峰值显存和吞吐，由郝天乐据此估算训练、生成及重跑时间，三人共同更新排期。

\textbf{11 月 23 日至 12 月 6 日：Caption 验收与中期原型。}单卡先完成 A/B 训练和统一四指标评价，再运行 B-global 消融；按验证 CIDEr-D 选定 A 权重。罗开炜并行准备迁移接口，取得正式权重后完成 Hybrid 原型，郝天乐复跑真实样本。12 月 3--5 日联调并交叉复算，张嘉荣整理同图案例与中期资料，6 日提交。中期验收要求真实数据链路、可加载权重、四指标及失败统计均可出示，不以 tiny smoke 替代正式原型记录。

\textbf{12 月 7--20 日：迁移、适配与奖励优化。}7--10 日完成 Hybrid 训练及两种迁移对照；11--16 日由郝天乐运行 Qwen 同规模原始/微调配对实验，罗开炜复核监督范围，张嘉荣在非测试样本上完善 Judge 和人评工具。17--20 日由罗开炜完成 SCST 与等步数继续交叉熵对照，郝天乐检查评价覆盖与训练记录。各模型均依据验证集选择；20 日统一冻结权重、数据版本、解码和评价协议，同时锁定人评样本 ID 与规则。

\textbf{12 月 21--28 日：正式评价、人工复核与展示。}21 日起生成冻结版本的测试预测和同图匿名候选，随后运行 Judge；22--24 日郝天乐与张嘉荣独立盲评，罗开炜复核分歧，张嘉荣汇总统计。按第五章约 60 张图像、A/B/A+SCST 三组输出估算，每名评分者需评价约 180 条描述，每条含四维评分及简短依据。25 日验收结果表、人评记录和可运行统计，26--27 日用于差异复核、展示整合与演练，28 日预留口头汇报。

\textbf{12 月 29 日至 1 月 17 日：分析与交付。}开展分层误差分析和独立复现，整理 Notebook/Markdown 报告及展示材料，1 月 14 日完成组内交付检查，15--16 日修正材料和复现说明，17 日统一提交。冻结后的测试结果仅用于评价与分析，不据此调整模型、提示或后处理；运行故障可按相同冻结配置重跑，并保留原因与日志。

\begin{table}[!ht]
\caption{共同里程碑与验收后交付}
\label{tab:project-schedule}\centering\small
\setlength{\tabcolsep}{4pt}\renewcommand{\arraystretch}{1.12}
\begin{tabularx}{\linewidth}{@{}P{2.0cm}P{1.65cm}Y@{}}\toprule
\textbf{日期} & \textbf{牵头人} & \textbf{交付条件}\\\midrule
11 月 15 日 & 张嘉荣 & 开题报告、研究问题、模型对照及个人检查点。\\
11 月 22 日 & 郝天乐 & 可追溯数据清单、指标预检查、GPU 试跑与修订排期。\\
12 月 6 日 & 郝天乐 & A/B 四指标、选定 A 权重、Hybrid 原型及中期资料。\\
12 月 20 日 & 罗开炜 & 计划实验的验证记录与选定权重，冻结测试和人评协议。\\
12 月 25 日 & 张嘉荣 & 统一结果表、Judge 与双人原始评分、可复算分析及展示版本。\\
1 月 14 日 & 罗开炜 & 跨成员复现通过，报告、展示材料与大文件获取说明齐备。\\
1 月 17 日 & 罗开炜 & Notebook/Markdown 报告，另交幻灯片或 2--5 分钟视频。\\\bottomrule
\end{tabularx}
\end{table}

\textbf{缓冲与延期处理。}11 月 22 日、12 月 6 日及 16 日检查关键路径。若数据映射或单卡吞吐不能支持原排期，由组长在下一次训练前更新作业顺序：保障两项必选及已选 SCST、Judge 的完整交付，优先调整额外种子、扩展消融与图表美化；保留约定的核心对照和失败统计。数据未锁定时继续接口及任务二数据核验，不启动未经审计的 Caption 正式训练。若核心实验仍延期，同步更新冻结和人评窗口，至少预留三天双人评分与一天统计复核；汇报如实区分已完成和待完成内容，不将未执行实验写成结果。

三次课程材料分别于 11 月 15 日、12 月 6 日及次年 1 月 17 日 24:00 前由组长提交。中期选择书面报告或口头报告；中期与结题口头汇报暂按 12 月 7 日和 12 月 28 日预留，具体形式及日期以教师最新通知为准。

